% GENERATED FILE. Edit canonical lessons, then run npm run generate:books.
% canonical-source-hash: fnv1a64:96d32044e57869bb
% canonical-lessons: FR-C68-banane, FR-C68-orange, FR-C68-terre, FR-C68-tomate, FR-C68-oignon

\chapter{Banana, Orange, Potato, Tomato, Onion}
\label{ch:fr-banana-orange-potato-tomato-onion}
\hlchaptermodality{68}

\begin{chapteropening}
\textbf{By the end of this chapter:} \emph{I can name a banana, an orange, a potato, a tomato, an onion.}

Complete the last lesson of chapter 68: I can name a banana, an orange, a potato, a tomato, an onion.
\end{chapteropening}

\section[la banane]{la banane — a banana}

\label{lesson:FR-C68-banane}



\begin{quote}
Before the new one: say the French for a fruit, then the French for an apple.
\end{quote}

\subsection*{You'll want to know: la banane}
\textbf{la banane} — \textquotedblleft{}a banana\textquotedblright{}.

\begin{cousinweb}
\textbf{banane} $\leftarrow$ Portuguese \emph{banana}, from a West African language, probably Wolof — borrowed. (EtymDB 2.1)
\end{cousinweb}

\begin{grammarlens}[title={What you've built}]
One more word for this chapter.
\end{grammarlens}

\subsection*{Your turn}
\begin{itemize}
\raggedright
  \item \emph{Say it:} \emph{la banane}
  \item \emph{Say it:} \emph{la banane}, once more
  \item \emph{Say it:} say \emph{la banane}
  \item \emph{Recall:} say the French for a fruit, then the French for an apple, then say \emph{la banane} again
\end{itemize}

\begin{tcolorbox}[breakable,colback=teal!4,colframe=teal!35!black,title={Before you move on}]
What does la banane mean? (\textquotedblleft{}A banana\textquotedblright{}.) Say it once more.
\end{tcolorbox}
\section[l'orange]{l'orange — an orange}

\label{lesson:FR-C68-orange}



\begin{quote}
Before the new one: say the French for an apple, then the French for a banana.
\end{quote}

\subsection*{You'll want to know: l'orange}
\textbf{l'orange} — \textquotedblleft{}an orange\textquotedblright{}.

\begin{cousinweb}
\textbf{orange} $\leftarrow$ Old French \emph{pome orenge}, \textquotedblleft{}orange apple\textquotedblright{} (EtymDB 2.1), which came through Italian from Arabic \emph{naranj} and Persian \emph{narang} — borrowed. The \emph{n-} fell away on the journey.
\end{cousinweb}

\begin{grammarlens}[title={What you've built}]
One more word for this chapter.
\end{grammarlens}

\subsection*{Your turn}
\begin{itemize}
\raggedright
  \item \emph{Say it:} \emph{l'orange}
  \item \emph{Say it:} \emph{l'orange}, once more
  \item \emph{Say it:} say \emph{l'orange}
  \item \emph{Recall:} say the French for an apple, then the French for a banana, then say \emph{l'orange} again
\end{itemize}

\begin{tcolorbox}[breakable,colback=teal!4,colframe=teal!35!black,title={Before you move on}]
What does l'orange mean? (\textquotedblleft{}An orange\textquotedblright{}.) Say it once more.
\end{tcolorbox}
\section[la pomme de terre]{la pomme de terre — a potato}

\label{lesson:FR-C68-terre}



\begin{quote}
Before the new one: say the French for a banana, then the French for an orange.
\end{quote}

\subsection*{You'll want to know: la pomme de terre}
\textbf{la pomme de terre} — \textquotedblleft{}a potato\textquotedblright{}. Literally \textquotedblleft{}apple of the earth\textquotedblright{}.

\begin{cousinweb}
\textbf{pomme de terre} — \textquotedblleft{}apple of the earth\textquotedblright{}: \emph{pomme}, \textquotedblleft{}apple\textquotedblright{}, \emph{de}, \textquotedblleft{}of\textquotedblright{}, and \emph{terre}, \textquotedblleft{}earth\textquotedblright{}, $\leftarrow$ Latin \emph{terra}. (EtymDB 2.1)
\end{cousinweb}

\begin{grammarlens}[title={What you've built}]
One more word for this chapter.
\end{grammarlens}

\subsection*{Your turn}
\begin{itemize}
\raggedright
  \item \emph{Say it:} \emph{la pomme de terre}
  \item \emph{Say it:} \emph{la pomme de terre}, once more
  \item \emph{Say it:} say \emph{la pomme de terre}
  \item \emph{Recall:} say the French for a banana, then the French for an orange, then say \emph{la pomme de terre} again
\end{itemize}

\begin{tcolorbox}[breakable,colback=teal!4,colframe=teal!35!black,title={Before you move on}]
What does la pomme de terre mean? (\textquotedblleft{}A potato\textquotedblright{}.) Say it once more.
\end{tcolorbox}
\section[la tomate]{la tomate — a tomato}

\label{lesson:FR-C68-tomate}



\begin{quote}
Before the new one: say the French for an orange, then the French for a potato.
\end{quote}

\subsection*{You'll want to know: la tomate}
\textbf{la tomate} — \textquotedblleft{}a tomato\textquotedblright{}.

\begin{cousinweb}
\textbf{tomate} $\leftarrow$ Spanish \emph{tomate} $\leftarrow$ Nahuatl \emph{tomatl} — borrowed. (EtymDB 2.1)
\end{cousinweb}

\begin{grammarlens}[title={What you've built}]
One more word for this chapter.
\end{grammarlens}

\subsection*{Your turn}
\begin{itemize}
\raggedright
  \item \emph{Say it:} \emph{la tomate}
  \item \emph{Say it:} \emph{la tomate}, once more
  \item \emph{Say it:} say \emph{la tomate}
  \item \emph{Recall:} say the French for an orange, then the French for a potato, then say \emph{la tomate} again
\end{itemize}

\begin{tcolorbox}[breakable,colback=teal!4,colframe=teal!35!black,title={Before you move on}]
What does la tomate mean? (\textquotedblleft{}A tomato\textquotedblright{}.) Say it once more.
\end{tcolorbox}
\section[l'oignon]{l'oignon — an onion}

\label{lesson:FR-C68-oignon}



\begin{quote}
Before the new one: say the French for a potato, then the French for a tomato.
\end{quote}

\subsection*{You'll want to know: l'oignon}
\textbf{l'oignon} — \textquotedblleft{}an onion\textquotedblright{}.

\begin{cousinweb}
\textbf{oignon} $\leftarrow$ Old French \emph{oingnon} $\leftarrow$ Latin \emph{unio}, \emph{unionis}, a kind of onion that grows as a single bulb — inherited. (EtymDB 2.1) \emph{unio} is built on \emph{unus}, \textquotedblleft{}one\textquotedblright{} (Lewis \& Short).
\end{cousinweb}

\begin{grammarlens}[title={What you've built}]
One more word for this chapter.
\end{grammarlens}

\subsection*{Your turn}
\begin{itemize}
\raggedright
  \item \emph{Say it:} \emph{l'oignon}
  \item \emph{Say it:} \emph{l'oignon}, once more
  \item \emph{Say it:} say \emph{l'oignon}
  \item \emph{Recall:} say the French for a potato, then the French for a tomato, then say \emph{l'oignon} again
\end{itemize}

\begin{tcolorbox}[breakable,colback=teal!4,colframe=teal!35!black,title={Before you move on}]
What does l'oignon mean? (\textquotedblleft{}An onion\textquotedblright{}.) Say it once more.
\end{tcolorbox}
